% Source: https://artofproblemsolving.com/community/c3h1845607_interesting_problem_on_number_theory
% Stable identity: aops-topic-1845607-q01
\textbf{$z\ge0$.} Then $N=2^z+2$ is an integer. If $N=\left(\dfrac {p}{q}\right)^2$ in lowest terms, then $q^2N=p^2$; a prime factor of $q$ would divide $p$, so $q=1$ and $N$ is a perfect square. $z=0$ gives $3$, not a square. $z=1$ gives $4=2^2$. For $z\ge2$, $N\equiv2\pmod4$, while squares are $\equiv0$ or $1\pmod4$. So only $z=1$.

\textbf{$z=-m$ with $m\ge1$.} Then
\[
2^{-m}+2=\dfrac{1+2^{m+1}}{2^m},
\]
and this fraction is in lowest terms because its numerator is odd. If it equals $\left(\dfrac {p}{q}\right)^2=\dfrac{p^2}{q^2}$ with $\dfrac {p}{q}$ in lowest terms ($p,q>0$), then $\dfrac{p^2}{q^2}$ is also in lowest terms. Since a fraction has only one representation in lowest terms with a positive denominator, $q^2=2^m$ and $p^2=1+2^{m+1}$. The second equation gives $2^{m+1}=(p-1)(p+1)$. So $p-1$ and $p+1$ are powers of $2$ differing by $2$, namely $2$ and $4$. Hence $p=3$ and $m+1=3$, so $m=2$; then $q^2=4$ works. So $z=-2$, where $2^{-2}+2=\dfrac{9}{4}=\left(\dfrac{3}{2}\right)^2$.

The integers are $\boxed{z=1\text{ and }z=-2}$.
